\centerline{\bf \Huge Теорема Виета}

1. Можно ли подобрать три различных действительных числа так, чтобы их сумма была равна числу $a$, сумма попарных произведений была равна числу $a^2$, а произведение равно числу $a^3$ ?

\textbf{Решение.} Предположим, что такие числа существуют. Обозначим их через $x, y$ и z. Тогда

$$
x+y+z=a, \quad x y+y z+z x=a^2, \quad x y z=a^3 .
$$


Заметим, что по теореме Виета для кубического многочлена числа $x, y, z$ являются корнями многочлена $f(t)=t^3-a t^2+a^2 t-a^3=\left(t^2+a^2\right)(t-a)$. Но у такого многочлена есть не более двух различных корней - противоречие.

Ответ. Нет.

\problem{2}
Вещественные числа $a, b$ и $c$ таковы, что $a+b+c>0, a b+b c+c a>0$ и $a b c>0$. Докажите, что числа $a, b$ и $c$ положительные.

\textbf{Решение.} Рассмотрим многочлен $f(t)=(t+a)(t+b)(t+c)$. Числа $-a,-b$ и $-c$ являются его корнями. Но у этого многочлена все коэффициенты положительны (они равны $1, a+b+c, a b+b c+c a$ и $a b c$ ), поэтому корни должны быть отрицательными. Значит, $-a,-b,-c<0$, что и требовалось доказать.

\problem{3}
Вещественные числа $a, b$ и $c$ таковы, что $\frac{1}{a}+\frac{1}{b}+\frac{1}{c}=\frac{1}{a+b+c}$. Докажите, что среди чисел $a, b, c$ найдутся два числа, дающие в сумме 0 .

\textbf{Решение.} Рассмотрим многочлен $f(t)=(t-a)(t-b)(t-c)$. Введем обозначения:

$$
a+b+c=u, \quad a b+b c+c a=v, \quad a b c=w .
$$


Тогда условие задачи записывается так: $\frac{v}{w}=\frac{1}{u}$, откуда $w=u v$. Преобразуем многочлен $f$ :

$$
f(t)=t^3-u t^2+v t-w=t^3-u t^2+v t-u v=\left(t^2+v\right)(t-u) .
$$


Раз у этого многочлена есть три вещественных корня, то $v<0$ и $\pm \sqrt{-v}$ - два из них. Их сумма и равна 0.

\problem{4}
Петя придумал действительные числа $a, b, c$ и сообщил про них Васе, что

$$
a^2+b^2+c^2+a b c=5 \quad \text { и } \quad a+b+c=3 .
$$


Может ли Вася назвать число, которое равно одному из придуманных Петей?
\textbf{Решение.} Введем обозначения:

$$
a+b+c=u, \quad a b+b c+c a=v, \quad a b c=w .
$$

Тогда условие задачи записывается так:

$$
u=3, \quad\left(u^2-2 v\right)+w=5 .
$$


Отсюда находим $w=2 v-4$. Рассмотрим многочлен $f(t)=(t-a)(t-b)(t-c)$. Тогда

$$
f(t)=t^3-u 6^2+v t-w=t^3-3 t^2+v t-2 v+4 .
$$


Заметим, что $f(2)=0$, поэтому одно из чисел $a, b$ или $c$ обязано равняться 2 . Значит, если Вася назовет число 2 , то точно не ошибется.

Ответ. Да.

\problem{5}
Пусть $f$ - приведенный целочисленный многочлен степени $n$ с ненулевыми коэффициентами. Известно, что у этого многочлена есть $n$ различных целых попарно взаимно простых корней. Докажите, что его коэффициенты $a_n$ и $a_{n-1}$ взаимно просты.

\textbf{Решение.} Пусть $f(x)=x^n+a_1 x^{n-1}+\ldots+a_{n-1} x+a_n$ - наш многочлен, а $x_1, x_2, \ldots$, $x_n$ - его корни. По теореме Виета имеем:

$$
x_1 x_2 \ldots x_n=(-1)^n a_n, \quad x_1 x_2 \ldots x_{n-1}+\ldots+x_2 x_2 \ldots x_n=(-1)^{n-1} a_{n-1} .
$$

Предположим, что коэффициенты $a_{n-1}$ и $a_n$ не взаимно просты и оба делятся на простое число $p$. Тогда $x_1 x_2 \ldots x_n \vdots p$, поэтому существует корень, делящийся на $p$. Без ограничения общности это $x_1$. Тогда в выражении $x_1 x_2 \ldots x_{n-1}+\ldots+x_2 x_2 \ldots x_n= (-1)^{n-1} a_{n-1}$ все слагаемые в левой части, кроме последнего, содержат множитель $x_1$ и потому также делятся на $p$. Кроме того, на $p$ делится правая часть этого равенства. Значит, последнее слагаемое $x_2 x_3 \ldots x_n$ также делится на $p$. Поэтому есть еще один корень, отличный от $x_1$, кратный $p$, что невозможно.

\problem{6}
Решить в целых числах уравнение $x^3+y^3+6 x y=8$.

\textbf{Решение.} Положим $u=x+y, v=x y$. Тогда условие задачи переписывается в виде $u\left(u^2-3 v\right)+6 v=8$, откуда

$$
(u-2)\left(u^2+2 u+4-3 v\right)=0 .
$$


Значит, либо $u=2$ и пары ( $n ; 2-n$ ) являются решениями для любого целого $n$, либо выполнено равенство $u^2+2 u+4=3 v$. Заметим, однако, что для чисел $u$ и $v$ выполнено неравенство $u^2 \geqslant 4 v$ (это можно понять, просто раскрыв скобки, или заметив, что $u^2-4 v$ - дискриминант квадратного уравнения с корнями $x$ и $y$ ). Тогда

$$
u^2+2 u+4=3 v \leqslant \frac{3}{4} u^2 \quad \Longrightarrow \quad(u / 2+2)^2 \leqslant 0 .
$$


Значит, $u=-4$ и $v=4$, откуда находим $x=y=-2$.
Ответ. $(-2 ;-2)$ или $(n ; 2-n)$, где $n \in \mathbb{Z}$.